\chapter{Conventions Used in This Book}
\label{ch:conventions}

\section*{One stack, and what ``true'' means}

Every behavioral claim is written against the reference stack in \Cref{app:refenv}:
Spark 3.5.2 and Iceberg format version 2.
Three rules follow from that.

\begin{itemize}
  \item A behavior is stated as fact only if it holds on that stack.
  \item When behavior changed across versions, a Note names the version (for
  example, that adaptive coalescing is on by default from Spark 3.2).
  \item Defaults are quoted as of Spark 3.5.2. A number --- a threshold, a size, a
  speedup --- always comes with a source, or with ``measured'' and the environment
  that measured it.
\end{itemize}

\section*{Measured, derived, and claimed}

The book distinguishes three kinds of statement, because they deserve different
trust.

\begin{description}[leftmargin=0pt, labelindent=0pt, itemsep=4pt]
  \item[Measured.] Produced by a lab script and shown with its output. These are
  the numbers and plans in listings and figures. They are reproducible on the lab
  setup; on your cluster, expect the same mechanism and a different magnitude.
  \item[Derived.] Follows from a formula or from the source code of Spark or
  Iceberg, and is checked against a measurement where one exists.
  \item[Claimed.] Stated from documentation or experience, but not tested here.
  These are marked ``claim --- not tested'' or ``claim --- verify on your
  cluster.''
\end{description}

\section*{Listings}

Code and output appear in listings with a language tag: \texttt{sql}, \texttt{scala},
or \texttt{text} for plans, UI excerpts, and shell output. Query plans are lightly
trimmed; an ellipsis comment such as \texttt{-- ... 3 Filter nodes elided} marks a
cut. Wide plans use a smaller type and wrap long lines rather than running off the
page. Configuration keys in running text are typeset as
\conf{spark.sql.adaptive.enabled}, and may break at the dots.

A listing that produces output is followed by that output, labelled, so you can
compare what you expect with what happened.

\section*{Terminology}

Each term is defined once, in italics, at its first use, and used plainly after
that. Where two names circulate for the same thing, the book picks one:

\begin{itemize}
  \item \emph{stage boundary}, not ``shuffle boundary'';
  \item \emph{unified memory}, not ``managed memory'';
  \item \emph{manifest list}, \emph{manifest}, and \emph{data file}, exactly as the
  Iceberg specification spells them;
  \item an \emph{executor} is a process, a \emph{task} is a thread, and a
  \emph{slot} is one core.
\end{itemize}

The glossary (\Cref{app:glossary}) lists every term with the chapter that defines
it.

\section*{Figures and cross-references}

Figures are vector graphics, so they stay sharp at any zoom, and share one color
palette. Figures that show a measured result say so in the caption. Cross-references are links; a reference to a
chapter, section, figure, or listing is always followed by its number.

\section*{Callouts}

Four callouts mark the text that is not part of the main line of the argument. At
most one appears on a typical page.

\begin{note}
Marks a version gotcha, a clarification, or a pointer to a later chapter.
\end{note}

\begin{warning}
Marks a default or a pattern that silently costs money, correctness, or data.
\end{warning}

\begin{tip}
Marks a diagnostic shortcut --- a UI panel to check, a one-line command.
\end{tip}

\begin{analogy}
Marks a plain-language comparison used to build intuition for a mechanism ---
never a substitute for the mechanism itself, and always pushed only as far as it
stays accurate. Skippable if it doesn't click for you; the surrounding prose
never depends on it.
\end{analogy}

\section*{The running example}

The \texttt{orders} table and its star schema are introduced in \Cref{app:refenv}
and reused in every chapter. Table and column names are typeset in
\texttt{monospace}.
